\documentclass[10pt,a4paper]{amsart}
\usepackage[margin=19mm]{geometry}
\usepackage{amsmath,amssymb,mathtools}
\usepackage[scr=rsfs,cal=euler]{mathalfa}
\usepackage{xfrac}
\usepackage[unicode,hidelinks,bookmarks=false]{hyperref}
\newcommand{\ie}{i.e.\ }
\newcommand{\cf}{cf.\ }
\newcommand{\resp}{respectively\ }
\newcommand{\Subsection}{\subsection}
\providecommand{\nobreakdash}{\nobreak\hbox{-}}
\setlength{\emergencystretch}{2em}
\multlinegap0pt
\usepackage{bm}
\usepackage{bbm}
\usepackage{dsfont}
\usepackage{xspace}
\usepackage{cancel}
\usepackage{mathtools}
\usepackage{amsthm}
\makeatletter
\@ifundefined{thm}{\newtheorem{thm}{Theorem}}{}
\makeatother
\title[An ${\mathfrak S}\_3$-cover of $K\_4$ and integral polyhedral ]{An ${\mathfrak S}\_3$-cover of $K\_4$ and integral polyhedral graphs}
\author{Taizo Sadahiro}
\date{Source: arXiv:2508.18593v2 (CC BY 4.0)}
\begin{document}
\maketitle

\noindent\emph{Source and licence.} An ${\mathfrak S}_3$-cover of $K_4$ and integral polyhedral graphs. Version arXiv:2508.18593v2 (CC BY 4.0), distributed under the Creative Commons Attribution 4.0 International licence, as recorded in the arXiv licence metadata for this exact version and shown on the abstract page https://arxiv.org/abs/2508.18593v2. Full rights notice: licenses/arxiv-2508.18593-CC-BY-4.0.txt.\par\smallskip

\noindent\textit{Editorial scope.} Counted unit: the Thm whose source label is \texttt{thm:zeta-relation-S3} in the article (source0.tex lines 748--759, source label \texttt{thm:zeta-relation-S3}), delivered together with the article's own complete original proof of it (source0.tex lines 761--823, one \texttt{proof} environment of 63 lines). No mathematics is written, repaired or re-derived: statement and proof are the article's own text line for line. Required context retained uncounted: eq:Lfun1 (lines 711--714); eq:Linducedrep (lines 722--725); eq:zetafactorization (lines 729--735). Every definition, display and label of those lines is retained unchanged. Omitted from this package and not redistributed: 1--710, 715--721, 726--728, 736--747, 760--760, 824--1342. Nothing inside a retained excerpt is omitted or altered: every retained body is delivered exactly as the article's lines are written, so no omission receipt of any kind accompanies this package. Citations: the retained excerpts carry no citation command at all, so no bibliography entry of the article is transcribed and no reference is redistributed. Reference adaptation: none was needed and none was made; every reference inside the delivered unit resolves inside it, so no unresolved reference or citation is produced. Printed slips kept literal: no printed word, symbol, hypothesis or number of the counted statement or of its proof was altered. Layout and class: the article's own class is \texttt{article} with packages including amsmath, amssymb, amsthm, tikz, pgf, float, tikz-3dplot; the wrapper is the lane's pinned \texttt{amsart} editorial wrapper at 10pt on A4, which restates the theorem-like environments the retained text needs and adds the article's own macro definitions that the retained text needs (none), with extra wrapper packages (bm, bbm, dsfont, xspace, cancel, mathtools). The delivered numbering is the unit's own. Assets: no retained span contains a figure environment, a graphics-inclusion command, a PDF-page inclusion, a table or a TikZ picture; no image file is copied into this package, the package declares no asset dependency and no image file is redistributed with it. Licence: version arXiv:2508.18593v2 of this article is distributed under the Creative Commons Attribution 4.0 International licence, as recorded in the arXiv licence metadata for this exact version and shown on the abstract page https://arxiv.org/abs/2508.18593v2. A full rights notice is delivered at licenses/arxiv-2508.18593-CC-BY-4.0.txt. Characters: every retained excerpt is pure ASCII, so the delivered engine, XeLaTeX with its default Latin Modern text font, reports no missing character.
 \begin{equation}
  \label{eq:Lfun1}
  L(u,\rho,\Gamma/\Delta) =  L(u, \widetilde{\rho}, \widetilde{\Delta}/\Delta).
 \end{equation}

 \begin{equation}
  \label{eq:Linducedrep}
  L(u, \rho, \Gamma/M) = L(u, \rho^\#, \Gamma/\Delta).  
 \end{equation}

 \begin{equation}
  \label{eq:zetafactorization}
  \zeta_\Gamma(u)
 =
 \prod_{\chi\in \widehat{{\rm Gal}(\Gamma/\Delta)}}
 L(u, \chi, \Gamma/\Delta)^{d_\chi},
 \end{equation}

\begin{thm}
\label{thm:zeta-relation-S3}
Let $p:Y\to X$ be a Galois cover of finite graphs with
$\mathrm{Gal}(Y/X)\cong {\mathfrak S}_3$.  Let
$Q=Y/\langle(1,2,3)\rangle$ and $T=Y/\langle(1,2)\rangle$ be the
intermediate coverings corresponding to the cyclic subgroups 
$C_3=\langle(1,2,3)$ and $C_2=\langle(1,2)\rangle$ respectively.  
Then we have the identity
\[
\zeta_Y(u)\,\zeta_X(u)^2 = \zeta_Q(u)\,\zeta_T(u)^2.
\]
\end{thm}

\begin{proof}
%Write $L(u,\rho,Y/X)$ for the Artin $L$-function associated with an
%irreducible representation $\rho$ of ${\mathfrak S}_3$ (for the covering $Y\to X$), and 
Write $L(u,\chi,Y/H)$ for the $L$-function of a representation $\chi$ of a
subgroup $H$ (for the covering $Y\to Y/H$).  
Since ${\mathfrak S}_3$ has three  irreducible representations,
the trivial representation $1$, the sign representation ${\rm sgn}$,
and the standard representation ${\rm std}$,
the factorization formula  $(\ref{eq:zetafactorization})$ yields
\begin{equation}
 \label{eq:factorZetaY}
\zeta_Y(u)=\zeta_X(u)\,L(u,\mathrm{sgn},Y/X)\,L(u,\mathrm{std},Y/X)^2.
\end{equation}

Since the subgroup $C_3=\langle(1\,2\,3)\rangle$ is normal in ${\mathfrak S}_3$,
$Q=Y/C_3$ is a Galois cover of $X$ whose Galois group
is ${\mathfrak S}_3/\langle(1,2,3)\rangle\cong C_2$, we have a factorization
\[
 \zeta_Q(u) = \zeta_X(u)L(u, {\rm sgn}_Q, Q/X),
\]
where ${\rm sgn}_Q$ is the sign representation of $Gal(Q/X)$.
By $(\ref{eq:Lfun1})$, we have
%$\mathrm{Ind}_{C_3}^{S_3}1 \cong 1\oplus\mathrm{sgn}$, hence
\[
L(u, {\rm sgn}, Y/X) = L(u, {\rm sgn}_Q, Q/X),
%\zeta_Q(u)=L(u,1,Y/X)\,L(u,\mathrm{sgn},Y/X)=\zeta_X(u)\,L(u,\mathrm{sgn},Y/X),
\]
therefore we have
\begin{equation}
 \label{eq:Lsgn}
L(u,\mathrm{sgn},Y/X)=\frac{\zeta_Q(u)}{\zeta_X(u)}.
\end{equation}

For the subgroup $C_2=\langle(1\,2)\rangle$ note that
$\mathrm{Ind}_{C_2}^{{\mathfrak S}_3}\mathrm{sgn}_{C_2}\cong \mathrm{sgn}\oplus\mathrm{std}$
where ${\rm sgn}_{C_2}$ is the sign representation of $C_2$.
Then, by the induction property $(\ref{eq:Linducedrep})$
of $L$-functions, we have
\[
L(u,\mathrm{sgn},Y/T)=L\big(u,\mathrm{Ind}_{C_2}^{{\mathfrak S}_3}\mathrm{sgn},Y/X\big)
= L(u,\mathrm{sgn},Y/X)\,L(u,\mathrm{std},Y/X).
\]
Applying the factorization property $(\ref{eq:zetafactorization})$ to
the cover $Y\to Y/C_2$ whose Galois group is $C_2$, we have
\begin{eqnarray*}
 \zeta_Y(u) & = & \zeta_T(u) L(u,{\rm sgn}_{C_2}, Y/T)\\
 & = & \zeta_T(u) L(u,{\rm Ind}_{C_2}^{{\mathfrak S}_3}{\rm sgn}_{C_2}, Y/X)\\
 & =& \zeta_T(u)L(u,\mathrm{sgn},Y/X)\,L(u,\mathrm{std},Y/X)\\
 & = & \zeta_T(u)\frac{\zeta_Q(u)}{\zeta_X(u)}L(u,\mathrm{std},Y/X).
\end{eqnarray*}
Thus we have obtained
\begin{equation}
 \label{eq:Lstd}
  L(u,\mathrm{std},Y/X)
  = \frac{\zeta_Y(u)\,\zeta_X(u)}{\zeta_T(u)\,\zeta_Q(u)}.
\end{equation}
Substituting $(\ref{eq:Lsgn})$ and $(\ref{eq:Lstd})$ into
$(\ref{eq:factorZetaY})$ yields the claimed identity
\[
\zeta_Y(u)\,\zeta_X(u)^2 = \zeta_Q(u)\,\zeta_T(u)^2.
\]
This completes the proof.
\end{proof}

\end{document}
